\clearpage
\subsection{Logarithmen als Umkehrung des Exponenzierens}
\label{subsec:real_functions_logarithms}
Für $a>0$, $a\ne1$ ist $a^x$ streng monoton und nimmt jeden positiven
Wert genau einmal an. Deshalb lässt sich aus einem positiven
Funktionswert der zugehörige Exponent eindeutig zurückgewinnen.
\par\noindent\begin{minipage}{\linewidth}
\begin{DefinitionBox}{Logarithmus}{real_functions_logarithm}
  Für $a>0$, $a\ne1$ und $y>0$ ist der \emph{Logarithmus}
  \EnglishTerm{logarithm} $\log_a y$ der Exponent, mit dem die Basis $a$
  potenziert werden muss, um $y$ zu erhalten:
  \[
    a^x=y\quad\Longleftrightarrow\quad x=\log_a y.
  \]
  Die Logarithmusfunktion $\log_a:(0,\infty)\to\R$ ist die
  Umkehrfunktion von $x\mapsto a^x$.
\end{DefinitionBox}
\begin{ExampleBox}[before skip=0.3\baselineskip]{Gesucht ist der Exponent}{real_functions_logarithm_values}
  Übersetzen Sie jede Logarithmusfrage in eine Potenzgleichung.
  \[
    \begin{array}{c@{\qquad}c@{\qquad}c}
      2^{\,?}=8&10^{\,?}=0{,}01&(1/2)^{\,?}=4\\[0.5em]
      \log_2 8=\Cloze{1cm}{3}&
      \log_{10}0{,}01=\Cloze{1cm}{-2}&
      \log_{1/2}4=\Cloze{1cm}{-2}
    \end{array}
  \]
\end{ExampleBox}
\end{minipage}

\begin{center}
  \CachedTikzFigure{content/resources/chapters/real_functions/tikz/functions_logarithms.tex}
\end{center}

\begin{TheoremBox}{Eigenschaften der Logarithmusfunktion}{real_functions_log_properties}
  Die Funktion $\log_a x$ ist auf $(0,\infty)$ stetig und hat die
  \textcolor{green!55!black}{Wertemenge} $\R$. Es gilt $\log_a1=0$ und $\log_a a=1$.
  \begin{center}
    \renewcommand{\arraystretch}{1.4}
    \begin{tabular}{@{}lcc@{}}
      \toprule
       & $a>1$ & $0<a<1$\\ \midrule
      Monotonie & streng steigend & streng fallend\\
      $x\to0^+$ & $\log_a x\to-\infty$ & $\log_a x\to+\infty$\\
      $x\to+\infty$ & $\log_a x\to+\infty$ & $\log_a x\to-\infty$\\
      \bottomrule
    \end{tabular}
  \end{center}
  Die Gerade $x=0$ ist eine senkrechte Asymptote. Es gibt keinen
  $y$-Achsenabschnitt.
\end{TheoremBox}
Die besonderen Bezeichnungen
\emph{natürlicher Logarithmus} \EnglishTerm{natural logarithm}
$\ln x=\log_e x$ und \emph{dekadischer Logarithmus}
\EnglishTerm{common logarithm} $\lg x=\log_{10}x$ verwenden wir im Folgenden.

\begin{TheoremBox}{Logarithmengesetze}{real_functions_log_laws}
  Für $u,v>0$, $r\in\R$ sowie $a,b>0$ mit $a,b\ne1$ gelten:
  \[
    \begin{aligned}
      \log_a(uv)&=\log_a u+\log_a v,\\[0.6em]
      \log_a(u/v)&=\log_a u-\log_a v,\\[0.6em]
      \log_a(u^r)&=r\log_a u,\\[0.6em]
      \log_a u&=\frac{\log_b u}{\log_b a}
               =\frac{\ln u}{\ln a}.
    \end{aligned}
  \]
  Insbesondere gelten $\ln(e^x)=x$ für $x\in\R$ und
  $e^{\ln x}=x$ für $x>0$.
  Jede Exponentialfunktion lässt sich damit zur Basis $e$ schreiben:
  $a^x=e^{x\ln a}$ für $a>0$. Bei $0<a<1$ ist $\ln a<0$,
  entsprechend einer exponentiellen Abnahme.
\end{TheoremBox}
Die Produktregel folgt aus den Potenzgesetzen:
Mit $s=\log_a u$ und $t=\log_a v$ ist
$uv=a^s a^t=a^{s+t}$; der gesuchte Exponent ist also $s+t$.
\begin{RemarkBox}
  Für Summen gibt es keine entsprechende Zerlegung:
  $\ln(u+v)$ ist im Allgemeinen nicht $\ln u+\ln v$.
  Ein Logarithmus darf im Reellen nur auf ein positives Argument
  angewendet werden.
\end{RemarkBox}

\begin{ExampleBox}{Wann ist die Amplitude halbiert?}{real_functions_half_life}
  Für $A(t)=A_0e^{-t/\tau}$ mit $A_0>0$ und $\tau>0$ suchen wir die
  \emph{Halbwertszeit} \EnglishTerm{half-life} $t_{1/2}$.
  Berechnen Sie sie zunächst allgemein und anschließend für $\tau=2\,\mathrm{s}$.
  \begin{ClozeArea}[34]{11}
    Die Bedingung $A(t_{1/2})=A_0/2$ führt auf eine Gleichung für den Exponenten:
    \[
      \begin{array}{r@{\quad}l@{\qquad}l}
        &A_0e^{-t_{1/2}/\tau}=A_0/2&\TransformOp{:A_0}\\[0.9em]
        \Longleftrightarrow&e^{-t_{1/2}/\tau}=1/2&\TransformOp{\ln}\\[0.9em]
        \Longleftrightarrow&-t_{1/2}/\tau=\ln(1/2)&\TransformOp{\cdot(-\tau)}\\[0.9em]
        \Longleftrightarrow&t_{1/2}=\tau\ln2.&
      \end{array}
    \]
    Für $\tau=2\,\mathrm{s}$ ergibt sich $t_{1/2}\approx1{,}386\,\mathrm{s}$.
    Der Anfangswert $A_0$ beeinflusst die Halbwertszeit nicht.
  \end{ClozeArea}
\end{ExampleBox}
